\documentclass[10pt,a4paper]{amsart}
\usepackage{amsmath,amssymb,amsthm,graphicx}
\usepackage{verbatim}
\newenvironment{reference}{\comment}{\endcomment}
%\newenvironment{reference}{}{}
\newenvironment{slogan}{\comment}{\endcomment}
\newenvironment{history}{\comment}{\endcomment}

% For commutative diagrams we use Xy-pic
\usepackage[all]{xy}

% We use 2cell for 2-commutative diagrams.
\xyoption{2cell}
% Theorem environments.
%
\theoremstyle{plain}
\newtheorem{theorem}[subsection]{Theorem}
\newtheorem{proposition}[subsection]{Proposition}
\newtheorem{lemma}[subsection]{Lemma}

\theoremstyle{definition}
\newtheorem{definition}[subsection]{Definition}
\newtheorem{example}[subsection]{Example}
\newtheorem{exercise}[subsection]{Exercise}
\newtheorem{situation}[subsection]{Situation}

\theoremstyle{remark}
\newtheorem{remark}[subsection]{Remark}
\newtheorem{remarks}[subsection]{Remarks}

\numberwithin{equation}{subsection}

% Macros
%
\def\lim{\mathop{\mathrm{lim}}\nolimits}
\def\colim{\mathop{\mathrm{colim}}\nolimits}
\def\Spec{\mathop{\mathrm{Spec}}}
\def\Hom{\mathop{\mathrm{Hom}}\nolimits}
\def\Ext{\mathop{\mathrm{Ext}}\nolimits}
\def\SheafHom{\mathop{\mathcal{H}\!\mathit{om}}\nolimits}
\def\SheafExt{\mathop{\mathcal{E}\!\mathit{xt}}\nolimits}
\def\Sch{\mathit{Sch}}
\def\Mor{\mathop{\mathrm{Mor}}\nolimits}
\def\Ob{\mathop{\mathrm{Ob}}\nolimits}
\def\Sh{\mathop{\mathit{Sh}}\nolimits}
\def\NL{\mathop{N\!L}\nolimits}
\def\CH{\mathop{\mathrm{CH}}\nolimits}
\def\proetale{{pro\text{-}\acute{e}tale}}
\def\etale{{\acute{e}tale}}
\def\QCoh{\mathit{QCoh}}
\def\Ker{\mathop{\mathrm{Ker}}}
\def\Im{\mathop{\mathrm{Im}}}
\def\Coker{\mathop{\mathrm{Coker}}}
\def\Coim{\mathop{\mathrm{Coim}}}

% Boxtimes
%
\DeclareMathSymbol{\boxtimes}{\mathbin}{AMSa}{"02}

%
% Macros for moduli stacks/spaces
%
\def\QCohstack{\mathcal{QC}\!\mathit{oh}}
\def\Cohstack{\mathcal{C}\!\mathit{oh}}
\def\Spacesstack{\mathcal{S}\!\mathit{paces}}
\def\Quotfunctor{\mathrm{Quot}}
\def\Hilbfunctor{\mathrm{Hilb}}
\def\Curvesstack{\mathcal{C}\!\mathit{urves}}
\def\Polarizedstack{\mathcal{P}\!\mathit{olarized}}
\def\Complexesstack{\mathcal{C}\!\mathit{omplexes}}
% \Pic is the operator that assigns to X its picard group, usage \Pic(X)
% \Picardstack_{X/B} denotes the Picard stack of X over B
% \Picardfunctor_{X/B} denotes the Picard functor of X over B
\def\Pic{\mathop{\mathrm{Pic}}\nolimits}
\def\Picardstack{\mathcal{P}\!\mathit{ic}}
\def\Picardfunctor{\mathrm{Pic}}
\def\Deformationcategory{\mathcal{D}\!\mathit{ef}}
\title{Colimits of topological groups differing in topological spaces and topological groups}
\author{The Stacks Project Authors}
\date{stacks-project:examples.tex:local-snapshot-20261012; GNU FDL 1.2 or later}
\begin{document}
\maketitle
\noindent\textit{Editorial note.} All delivered mathematical text is the author's own, copied from the pinned Stacks Project chapter file examples.tex; no statement or proof step was rewritten, reordered or removed.\par\medskip
\noindent\textit{Editorial note.} The counted proof is the author's complete argument: The step 'it is easy to show that U cap G\_n is open' is the author's own terse carried claim (6895-6896), used verbatim; the irrationality of pi supplies the never-integral-multiple fact.\par\medskip
\noindent\textit{Editorial note.} The label lemma-colimit-topology names the counted statement, and the proof environment immediately following it in the source is delivered verbatim as its proof.\par\medskip
\noindent\textit{Editorial note.} The counted statement is carried by the author's own lemma environment in the chapter source.\par\medskip
\noindent\textit{Allowed standard prerequisites (editorial).}\begin{itemize}
\item Context span 6879-6887: G\_n = Q x R\textasciicircum{}n as a topological group with the Euclidean topology, closed embeddings G\_n -> G\_\{n+1\} appending a zero; G = colim G\_n with the colimit topology (U open iff G\_n cap U open for all n); the claim is that this is not a topological group
\item Context span 6889-6909: the set U = \{(x\_0, x\_1, ...) : |x\_j| < |cos(j x\_0)| for j > 0\}; since j x\_0 is never an integral multiple of pi/2 (pi irrational) it is easy to show U cap G\_n is open; if the colimit topology made G a topological group there would be an open neighbourhood V of 0 with V + V subset U; then for every j >= 0 some epsilon\_j > 0 has (0, ..., 0, x\_j, 0, ...) in V for |x\_j| < epsilon\_j; V + V subset U would force (x\_0, 0, ..., 0, x\_j, 0, ...) in U for |x\_0| < epsilon\_0 and |x\_j| < epsilon\_j; taking j with j epsilon\_0 > pi/2 and x\_0 in Q with |cos(j x\_0)| < epsilon\_j gives |cos(j x\_0)| < |x\_j| < epsilon\_j, contradicting membership in U
\end{itemize}\medskip
$G_n = \mathbf{Q} \times \mathbf{R}^n$, $n \geq 1$ seen as a topological group
for addition endowed with the usual (Euclidean) topology. Consider the closed
embeddings $G_n \to G_{n + 1}$ mapping $(x_0, \ldots, x_n)$ to
$(x_0, \ldots, x_n, 0)$. We claim that $G = \colim G_n$ endowed with the
topology
$$
U \subset G\text{ open} \Leftrightarrow G_n \cap U\text{ open }\forall n
$$
is not a topological group.

\medskip\noindent
To see this we consider the set
$$
U = \{(x_0, x_1, x_2, \ldots)\text{ such that }
|x_j| < |\cos(jx_0)| \text{ for } j > 0\}
$$
Using that $jx_0$ is never an integral multiple of $\pi/2$ as $\pi$
is not rational it is easy to show that $U \cap G_n$ is open. Since
$0 \in U$, if the topology above made $G$ into a topological group,
then there would be an open neighbourhood $V \subset G$ of $0$
such that $V + V \subset U$. Then, for every $j \geq 0$ there would
exist $\epsilon_j > 0$ such that $(0, \ldots, 0, x_j, 0, \ldots) \in V$
for $|x_j| < \epsilon_j$. Since $V + V \subset U$ we would have
$$
(x_0, 0, \ldots, 0, x_j, 0, \ldots) \in U
$$
for $|x_0| < \epsilon_0$ and $|x_j| < \epsilon_j$. However, if we
take $j$ large enough such that $j \epsilon_0 > \pi/2$, then we can
choose $x_0 \in \mathbf{Q}$ such that $|\cos(jx_0)|$ is smaller than
$\epsilon_j$, hence there exists an $x_j$ with
$|\cos(jx_0)| < |x_j| < \epsilon_j$. This contradiction proves the claim.

\begin{lemma}
\label{lemma-colimit-topology}
There exists a system $G_1 \to G_2 \to G_3 \to \ldots$ of (abelian)
topological groups such that $\colim G_n$ taken in the category of
topological spaces is different from $\colim G_n$ taken in the category
of topological groups.
\end{lemma}

\begin{proof}
See discussion above.
\end{proof}

\end{document}
